%======================================================================
%  Introduction
%
%  Source: CONVERSION_PROMPT.md section 8 (author-provided), snapshot
%          2026-08-28. Transcribed verbatim apart from LaTeX markup, the
%          \model macro, [ref] -> \cite, and the figure cross-references.
%
%  Two author-directed wording changes, both given in session 2026-08-28:
%    - paragraph 3: "...transfer to functional tasks and to sequences never
%      seen in training [ref, GPN; ref, Evo]. This ab initio generalization is
%      precisely what origin annotation on divergent plasmids demands."
%      -> "...transfer to applications like variant effect prediction,
%      annotation, and generation [ref, GPN; ref, Evo]."
%    - paragraph 4: "generalizes ab initio" -> "generalizes to unseen origins"
%      (standing instruction: avoid the term "ab initio").
%
%  2026-09-04: IMG/PR is now expanded on its first appearance in the body,
%  which is here. It was previously used unexpanded in the Abstract and twice
%  in this paragraph, and first expanded only in Results. The Abstract no
%  longer names it at all. This is the only change to the author-provided
%  Introduction text; the section is otherwise untouched and remains within
%  the Nature Biotechnology main-text budget.
%======================================================================

\section*{Introduction}
\label{sec:intro}

Plasmids are mobile, extrachromosomal DNA elements that shape bacterial
evolution and are central tools of microbial engineering. By transferring
genes horizontally between cells they spread traits --- antibiotic
resistance, virulence, and metabolic capabilities --- between otherwise
unrelated bacteria
% TODO cite: horizontal transfer of resistance/virulence/metabolic traits.
% The source text carries a bare [ref] here with no work named, and no
% candidate is present in any artifact bibliography.
, and by replicating independently of the chromosome they persist as stable
accessory genomes. The two processes that define a plasmid's life cycle also
make it useful in the laboratory: conjugative transfer moves engineered DNA
between hosts, and vegetative replication maintains it. Each is governed by a
short, non-coding cis-regulatory element --- the origin of transfer
(\emph{oriT}), where a relaxase nicks the plasmid to initiate conjugation,
and the origin of vegetative replication (\emph{oriV}), where a replication
initiator binds to begin replication (\fref{fig:pretraining}a,b).
Identifying these origins is therefore a prerequisite both for understanding
how a plasmid moves and replicates and for repurposing it as an engineering
vector.

Metagenomic sequencing has expanded the known plasmid universe by orders of
magnitude, cataloguing hundreds of thousands of plasmid sequences --- billions
of nucleotides --- most from uncultured lineages with no characterized host
\cite{camargo2024img}. Yet the elements that make these plasmids work remain
largely unannotated, and origins are among the hardest features to find.
\emph{oriT} and \emph{oriV} are short, intergenic, and structurally diverse,
defined less by conserved sequence motifs than by architecture --- iterons,
inverted repeats, and paired promoter boxes whose spacing and symmetry,
rather than their exact bases, carry function. The tools used to annotate
them (oriTfinder2 \cite{liu2025oritdb,li2018oritfinder}, OriV-Finder
\cite{li2025orivfinder}) are reference- and alignment-based: they recognize an origin by its
similarity to a curated set of experimentally characterized examples, so they
degrade sharply on the divergent origins that dominate metagenomic plasmids
and are effectively blind to origins unlike the reference set. The result is
a widening gap between the plasmids we can sequence and the ones we can
functionally interpret --- a gap that most constrains exactly the non-model,
uncultured plasmids of greatest interest for discovery and engineering.

Genomic language models (gLMs) offer a route past this reference dependence.
Trained by self-supervision to predict masked or successive nucleotides on
unannotated sequence, they learn representations that transfer to
applications like variant effect prediction, annotation, and generation
\cite{benegas2023dna,nguyen2024sequence}. General-purpose gLMs, however, are
trained across broad evolutionary scales and dominated by chromosomal
genomes; plasmids --- mobile, recombining, compositionally distinct, and
lacking any universal taxonomy --- are underrepresented, and pretraining on
the wrong sequence distribution can transfer worse than not pretraining at
all. Building a gLM for plasmids also raises a sampling problem absent for
cellular genomes: plasmids have no hierarchical classification to sample
across, and abundance in sequence databases reflects sequencing effort rather
than biological diversity, so naive training would over-fit a handful of
heavily sequenced clades.

Here we present \model, a genomic language model trained by masked-nucleotide
prediction on the plasmid sequences of the Integrated Microbial Genomes /
Plasmid Repository (IMG/PR). To learn from the
true diversity of the metagenomic plasmid universe rather than its sequencing
depth, we designed a diversity-aware sampling scheme that weights training
toward the breadth of plasmid diversity rather than sequencing depth
(\fref{fig:pretraining}). Using \model, we build a benchmark and a deployable
pipeline for plasmid-origin annotation: a lightweight classifier on the
model's frozen embeddings identifies origin-containing regions and, evaluated
on a homology-partitioned benchmark that isolates origins with no training
homolog, generalizes to unseen origins --- outperforming reference-based
tools and models trained from scratch, and matching genomic language models
up to 85$\times$ larger while using roughly 45$\times$ less inference
compute. Applied across all of IMG/PR, the classifier annotates putative
\emph{oriT} and \emph{oriV} genome-wide and reveals previously undescribed
transfer and replication gene architectures. Finally, interpreting the
trained model shows that it has internalized the cis-regulatory grammar of
these elements --- the repeats, iterons, and promoter couplings of
\emph{oriT} and \emph{oriV}, together with promoters and other functional
elements --- directly from sequence and without supervision. Together, these
results establish \model as a scalable, alignment-free foundation for
annotating and engineering the functional elements of diverse and uncultured
plasmids.
